\documentclass[letterpaper]{article}

\usepackage[letterpaper,margin=1in]{geometry}

% Compile from the repository root so this path resolves to the pinned WSP
% submodule.
\input{wsp/testing/test-case-library.tex}

\begin{document}
\def\TCID{TC-0001}
\def\TCTitle{Winsock member table}
\def\TCPurpose{Verify that the member table of the library describes the
\texttt{Winsock} class with named members of known kinds and events that are
methods.}
\def\TCPriority{\TCPriorityMedium}

\def\TCRequirementRef{REQ-0001}
\def\TCDesignRef{Not applicable.}
\def\TCFeatureRef{\texttt{include/wfc/mswinsck/mswinsck.hpp};
\texttt{src/mswinsck\_tables.cpp};
\texttt{tests/mswinsck\_tables\_test.cpp}.}
\def\TCTestTechnique{Structural inspection of the table: every class, member
and event is visited once.}

\def\TCPreconditions{The library is built (target \texttt{wfc\_mswinsck}) and
the test program \texttt{mswinsck\_tables\_test} exists.}
\def\TCEnvironment{Any supported target; the table is part of every build,
static or shared. CTest runs the program with a 120 second timeout.}
\def\TCAssumptions{Not applicable.}

\def\TCInputData{The table returned by \texttt{wfc::ui::mswinsck\_classes()}.}
\def\TCInitialState{None; the table is constant data.}

\def\TCProcedure{\begin{enumerate}
  \item Run \texttt{mswinsck\_tables\_test} (CTest \texttt{TC-0001-tables}, in WFC \texttt{wfc-mswinsck-TC-0001-tables}).
  \item The program reads the class list and compares its size and class
    names with the expected list, which is the single class
    \texttt{Winsock}.
  \item For each class it checks that the class has members, that each member
    has a non-empty name and a kind of \texttt{g}, \texttt{l}, \texttt{s} or
    \texttt{m}, and that each event has the kind \texttt{m}.
\end{enumerate}}
\def\TCExpectedResult{The program exits with status 0 and prints
\texttt{mswinsck: 1 classes}; no \texttt{FAIL} line is written.}
\def\TCPostConditions{None; the test writes no files.}

\TCTemplate

\end{document}